Constructing biomedical knowledge graphs from scientific literature traditionally requires domain-specific models that must be retrained for each therapeutic area, limiting generalizability across the diverse landscape of drug discovery research. Recent work on GraphRAG demonstrated that latent knowledge graphs can be extracted from unstructured text, but the resulting graphs encode similarity-based relationships derived from embedding proximity rather than scientifically grounded assertions. I present Epistract, an open-source Claude Code plugin that replaces similarity-based assembly with comprehension-based extraction. The system dispatches parallel LLM agents constrained by a biomedical domain schema comprising 17 entity types, 30 relation types, and over 40 ontology references. Extracted entities undergo molecular validation through RDKit for chemical structures and Biopython for protein and nucleic acid sequences, ensuring that graph nodes correspond to real molecular entities rather than text fragments. I validated the pipeline across six drug discovery domains---neurogenetics, oncology, rare disease, immuno-oncology, cardiovascular inflammation, and GLP-1 receptor agonist competitive intelligence---extracting 783 nodes and over 2,200 links from 111 documents with 100\% entity type coverage, 93\% relation type coverage, and all 25 user acceptance tests passing, with zero retraining between domains. The sixth scenario introduces multi-source corpus assembly from PubMed, Google Scholar, and Google Patents via SerpAPI, demonstrating extraction of molecular identifiers (peptide sequences, CAS numbers, InChIKeys) from patent documents---a workflow that mirrors how human researchers conduct competitive intelligence. The system is available as an MIT-licensed plugin installable in a single command, with all test corpora, extraction outputs, and graph evidence committed to the public repository for full reproducibility.\footnote{\url{https://github.com/usathyan/epistract}}
